\section{Constraint structure and the limits of reconstruction}
The gravitational branch now has a continuous account from transport geometry to constrained dynamics and spherical tests. The geometric examples identify the additional assumptions needed when a propagation rule is interpreted as gravity. Reduced Gowdy evolution and finite linearized ADM theory make constraint propagation concrete. The nonlinear product-rule obstruction then identifies a specific limitation of a finite pointwise discretization.

The alternatives address that limitation in their own settings. Matrix moment maps retain exact finite algebraic closure. Homological transfer requires a genuine fine differential and controlled contraction data. Resolved quadrature retains the stated polynomial identities under its degree budget, while nonlinear evolution can require additional spectral bands. These constructions therefore carry distinct state spaces and comparison problems.

The Lorentzian classical action supplies a gauge-theoretic parent. Its boundary description and local second-class reduction require the declared regularity, nondegeneracy, canonical transformation, and invertibility hypotheses. Projected discrete-gradient evolution preserves the stated autonomous invariants for an exact step solve within that chart. None of these hypotheses can be discarded when changing coordinates or moving toward a singular configuration.

Spherical mass and flux identities provide unusually direct checks. Electrovac benchmarks test a prescribed slicing trajectory of a fixed charged geometry. The scalar-collapse implementation adds evolved matter and radial constraint solves. Its trapped-sphere diagnostic is a bounded numerical result for the declared data and implementation. Establishing an event in a controlled continuum solution requires an additional existence and convergence argument.

\subsection{Mathematical work that remains}
A useful next nonlinear theorem should identify a compatible sequence of state spaces and maps, give a level-uniform stability estimate, and state the topology in which the constraint and evolution defects tend to zero. Exact evaluation at one resolution does not itself provide that sequence. A reduction theorem that crosses charts must also control degeneration of the gauge-fixing and Dirac matrices, boundary conditions, and the change of canonical variables.

The strong field route needs estimates adapted to its actual constraint solve and boundary closure. For scalar matter, this includes the nonlocal radial dependence of reconstructed mass and lapse, the branch of the slicing variables, and convergence strong enough to control the trapped-sphere diagnostic. Extending the construction to nonspherical geometry introduces gravitational degrees of freedom absent from a spherical model.

The inherited massive fermion conventions also need a physical comparison in flat space. Lorentzian signature, the Clifford relation, the kinetic factor, and the adjoint must be matched together. The algebraic gauge model does not settle that normalization.

A quantum theory requires further choices and arguments. The classical master equation does not define a quantum measure or settle anomalies. A common quantum geometry--matter construction and a controlled semiclassical comparison remain research obligations. The verification companion lists these obligations beside the source results they extend.

\subsection{From geometric to arithmetic refinement}
The next chapter develops register factorization, Haar packets, spectral reduction, and infinite-depth quantum hierarchies. It changes the model while preserving the discipline used here: finite identities, refinement compatibility, and analytic limits require separate arguments. The covariant fiber and cosmology chapters then examine whether a chosen refinement sector admits a physically controlled production history and consistent observable transport.

\begin{readerroadmap}[title={A concrete next comparison}]
Name the phase space, constraints, evolution, boundary conditions, and map between resolutions. State the regularity and chart on which the map is defined. Then prove the stability and convergence needed for the physical conclusion being sought.
\end{readerroadmap}
